\section*{Abstract}
\addcontentsline{toc}{section}{Abstract}

A \emph{copymint} takes an existing NFT image, edits it slightly, and mints it as new.
Kotzer et al.\ detect copymints with four perceptual image hashes (aHash, pHash, hsvHash
and sHash), each searched in its own BK-tree, and call an image a copy when at least two of
the four agree. Our project set out to make that detector cheaper to run. While
reimplementing it we found that sHash, the hash meant to survive cropping, does not define a
distance in the mathematical sense: comparing two images gives a different answer depending
on which image comes first, and the triangle inequality does not hold. A BK-tree depends on
both properties, so a search can silently skip a branch that holds a real copy.

We therefore replaced sHash with ORB feature matching (keypoints, descriptor matching and
RANSAC geometric verification) and kept everything else in the detector unchanged. On the
geometric edits sHash was built for, ORB reaches F1 90.6 against 76.4 for sHash (+14.2),
even though sHash was given a threshold chosen on the test set and ORB's was chosen on
training data. Held at equal precision (91.2\%), sHash catches 27.2\% of copies and ORB
90.0\%. In the full four-signal detector, with the paper's own thresholds left
untouched, swapping only sHash for ORB raises F1 from 77.5 to 86.4 (+8.9): recall rises
from 64.1\% to 77.2\% while precision stays at 98.1\%. When the paper's thresholds are
re-tuned to our data the gain is 81.7 to 87.5 (+5.8).

The improvement has a cost that we report as a finding. The paper stores its four hashes,
about 56 bytes, inside each mint transaction (300--500 bytes). ORB needs about 14.5\,KB per
image, roughly 36 transactions. Shrinking it removes the advantage before it fits: at 32
landmarks (1\,KB, still 2.6 transactions) ORB only ties sHash. In time, making a
fingerprint costs about the same for both (0.94\,ms for ORB, 1.46\,ms for sHash, both in
compiled code), but comparing one pair takes about 34\,ms for ORB against about
0.1\,\textmu s for sHash, so ORB needs an index such as LSH to scale.

Finally, we built an \emph{edit classifier} that predicts from the query image alone which
kind of edit was applied, so the detector could stop trusting signals known to fail on that
edit. The classifier is accurate (100\% on the ``detail destroyed'' decision, 97.5\% F1 on
``colour changed''), yet it does not improve the detector (F1 87.4 vs 87.5). A signal that
breaks goes silent rather than voting wrongly, so ignoring it changes nothing. A
multi-edit test then showed that the one situation where the classifier could help is not
observable with our cues: pixelation erases the colour trace the classifier relies on.
